\subsection{مبانی نظری و فرمول‌بندی ریاضی}
فرض کنید $\mathbf{W} \in \mathbb{R}^{d_{\text{out}} \times d_{\text{in}}}$ ماتریس وزن یک لایه خطی باشد. در الگوریتم \lr{ROAST}، ماتریس وزن به کاشی‌های مربعی با اندازه $B \times B$ افراز می‌شود:
\begin{equation}
T_{\text{out}} = \left\lceil \frac{d_{\text{out}}}{B} \right\rceil, \quad T_{\text{in}} = \left\lceil \frac{d_{\text{in}}}{B} \right\rceil
\end{equation}
تعداد کل کاشی‌ها برابر با $T_{\text{out}} \times T_{\text{in}}$ است. به جای ذخیره مقادیر مجزا، یک جدول وزن فشرده شامل $K$ باکت آموزش‌پذیر با پارامتر $\boldsymbol{\theta} \in \mathbb{R}^K$ تعبیه می‌شود ($K \ll d_{\text{out}} \cdot d_{\text{in}}$).

برای نگاشت هر کاشی در مختصات $(i, j)$ با $i \in \{0, \dots, T_{\text{out}}-1\}$ و $j \in \{0, \dots, T_{\text{in}}-1\}$ به یک باکت مشخص در جدول وزن، یک تابع درهم‌سازی با آنتروپی بالا و ترکیب بیتی (\lr{Bitwise Mixing}) فرمول‌بندی می‌شود:
\begin{align}
h_0 &= (i \cdot \mathtt{0x9E3779B1}) \oplus (j \cdot \mathtt{0x85EBCA77}) \\
h_1 &= h_0 \oplus (h_0 \gg 16) \\
h_2 &= h_1 \cdot \mathtt{0x7FEB352D} \\
h_3 &= (h_2 \oplus (h_2 \gg 15)) \ \& \ \mathtt{0x7FFFFFFF} \\
\operatorname{idx}(i, j) &= h_3 \bmod K
\end{align}
مقدار پارامتر کاشی $(i, j)$ از باکت استخراج می‌شود:
\begin{equation}
\mathbf{W}_{\text{tile}}(i, j) = \boldsymbol{\theta}[\operatorname{idx}(i, j)] \cdot \mathbf{1}_{B \times B}
\end{equation}
ماتریس کامل بازسازی‌شده با تکرار مقادیر کاشی‌ها تشکیل و به ابعاد مجاز برش داده می‌شود:
\begin{equation}
\hat{\mathbf{W}} = \left( \boldsymbol{\theta}[\operatorname{idx}(i, j)] \otimes \mathbf{1}_{B \times B} \right)_{:d_{\text{out}}, :d_{\text{in}}}
\end{equation}
همچنین در صورت استفاده از تطبیق کم‌مرتبه، اثر آداپتورهای \lr{LoRA} به صورت مستقیم در گذر پیشرو محاسبه می‌گردد:
\begin{equation}
\mathbf{y} = \mathbf{x} \hat{\mathbf{W}}^\top + \mathbf{b} + \frac{\alpha}{r} (\mathbf{x} \mathbf{A}_{\text{lora}}^\top) \mathbf{B}_{\text{lora}}^\top
\end{equation}


\paragraph{تحلیل طیفی و نامساوی‌های پایداری خطا}
برای بررسی میزان دقت تقریب ماتریس وزن یا نگاشت توجه در ماژول \lr{roast\_hashed\_linear}، تفاضل ماتریس فاکتورگیری‌شده از نگاشت متراکم اصلی بر اساس نامساوی هافدینگ-فون‌نیومان ارزیابی می‌گردد. فرض کنید $\mathbf{W} \in \mathbb{R}^{d_{\text{in}} \times d_{\text{out}}}$ ماتریس وزن هدف باشد که توسط تقریب کم‌مرتبه یا پراکنده با رتبه موثر $r$ جایگزین شده است. کران بالای خطای بازسازی بر مبنای نرم فروبنیوس به صورت زیر تصریح می‌شود:
\begin{equation}
\label{eq:spectral_bound_roast_hashed_linear}
\| \mathbf{W} - \mathbf{W}_{r} \|_F = \sqrt{\sum_{i=r+1}^{\min(d_{\text{in}}, d_{\text{out}})} \sigma_i^2(\mathbf{W})} \le \epsilon \| \mathbf{W} \|_F
\end{equation}
که در آن $\sigma_i(\mathbf{W})$ بیانگر $i$-امین مقدار منفرد مرتب‌شده ماتریس به صورت نزولی است. نسبت حفظ انرژی طیفی $\eta_r = \frac{\sum_{i=1}^r \sigma_i^2}{\sum_{i=1}^{d} \sigma_i^2}$ تضمین می‌کند که با انتخاب آستانه $\tau \ge 0.90$، بیش از نود درصد از پراکندگی واریانس داده‌ها حفظ شده و افت سیگنال در انتشار گرادیان رو به عقب به حداقل ممکن کاهش می‌یابد.

\paragraph{تحلیل مجانبی و پیچیدگی اشتراک حافظه سخت‌افزاری}
از منظر پیچیدگی محاسباتی و انتقال داده‌ها در سلسومراتب حافظه (حافظه سراسری \lr{HBM} به حافظه محلی \lr{SRAM})، عملیات ضرب ماتریسی در این مکانیزم نیازمند تخصیص بافر موقت با ابعاد فاکتورگیری‌شده است. برای پردازش دنباله‌ای به طول $N$ با بعد مخفی $D$، بار محاسباتی هسته از فرم درجه دو به رابطه زیر تقلیل می‌یابد:
\begin{equation}
\label{eq:complexity_roast_hashed_linear}
\mathcal{C}_{\text{roast\_hashed\_linear}} = \mathcal{O}\left( B \cdot N \cdot D \cdot r + B \cdot H \cdot N \cdot d_k \right) \ll \mathcal{O}\left( B \cdot H \cdot N^2 \cdot d_k \right)
\end{equation}
این بهبود مرتبه پیچیدگی، با کاهش چشمگیر حجم تبادل تانسورها در گذرگاه‌های ارتباطی کارت گرافیک، مانع از اشباع پهنای باند حافظه (پدیده \lr{Memory-Bound}) در فرآیند استنتاج طولانی‌مدت می‌گردد.
